% Section 3, part 1 of 8; label prefix: gau. The main file inputs only this part; each part inputs the next.
% Part 1: introduction; the cyclic form of PF4 (objects, statements, the conifold expansion).
% Part 2: the quintic system and the pulled-back model; Theorem U (notation, statement).
% Part 3: the commutant, lifting, the scalar equation.
% Part 4: proof of Theorem U, the leading term, the Z/5-degree, graded gauges.
% Part 5: (R1), the diagonal, the vanishing solutions.
% Part 6: the gauge identity (G), the canonical gauge, symplecticity, HPF.
% Part 7: polynomiality, formal WKB uniqueness, colour covariance of the model columns and of the actual rows.
% Part 8: the normalisation series D, the rows in the ring, the colours of D, (F_g).
% The proofs of the cyclic form of PF4 are in the appendix file a2_cyclicPF4. No new macros.
\section{Cyclic systems, Theorem U and the canonical gauge}
\label{sec:gau}

We compare the cyclic systems of the formal quintic at the conifold point and of the model: both have one simple pole and equal products (\cref{sec:gau-cyc,sec:gau-systems}), their gauges form a free module of rank one (Theorem~U, \cref{sec:gau-U}), (G) singles out the symplectic canonical gauge $\Pc$ (\cref{sec:gau-van}), and $\Pc$ applied to the model columns is the matrix of rows times the normalisation series $\Dn$ (\cref{sec:gau-R0}).

\emph{Conventions.} $\Dq=q\,\dd/\dd q=-(1-u)\,\dd/\dd u$ (\cref{def:set-genus0,def:set-points}); for a series in $u$ a prime denotes $\dd/\dd u$; $z$ is a constant, and $z\de_u$ is $z\,\dd/\dd u$ at fixed $z$, so $z\Dq=-(1-u)z\de_u$. $E_{ab}$ ($a,b\in\bbZ/5$) are the matrix units and $\mathbf e_a$ the standard basis vectors; solutions, gauges and rows are in lowered components (\cref{def:set-rows}).

%-------------------------------------------------------------------------------
\subsection{The cyclic form of the fourth-order operator}
\label{sec:gau-cyc}

In this subsection $\PFfour=\sum_{k=0}^4A_k\Dq^k$ is written in the variable $u$ (\cref{def:set-genus0}):
\[
A_4=u,\qquad A_3=-2(1-u),\qquad A_2=-\tfrac75(1-u),\qquad A_1=-\tfrac25(1-u),\qquad A_0=-\tfrac{24}{625}(1-u),
\]
and we use $\Dq(1-u)=1-u$ and $\Dq u=-(1-u)$ throughout. The proofs of \cref{lem:gau-confexp,thm:gau-cyc,cor:gau-cycsys} are in \cref{sec:gau-cycproof}.

\begin{defn}[Objects of the cyclic-form theorem]
\label{def:gau-conc}
\begin{enumerate}
\item $\bbK:=\bbQ\lser{u}$ with the derivation $\Dq$, and $\bbK[\ell]$ is the polynomial ring in an indeterminate $\ell$ with the unique extension of $\Dq$ for which $\Dq\ell=1-u^{-1}$ (so $\ell$ plays the role of $\log u$). $\bbK[\Dq]$ is the ring of differential operators $\sum_ka_k\Dq^k$, $a_k\in\bbK$, with $\Dq\circ a-a\circ\Dq=\Dq(a)$; it acts on every differential ring that extends $(\bbK,\Dq)$.
\item $f_0,f_1,f_2\in\bbQ\fps u$ are the solutions of $\PFfour$ with $f_k=u^k+O(u^3)$, and $\varpi_1=f_1+\gamma f_2$, $\gamma=\frac7{10}$, is the vanishing period (\cref{lem:set-frobbasis}). $y_3:=\varpi_1\ell+\rho_3\in\bbK[\ell]$, $\rho_3\in u^3\bbQ\fps u$, is the logarithmic solution of \cref{lem:set-frobbasis}~(iii), whose indeterminate $\ell$ is the one above; thus $\PFfour y_3=0$ in $\bbK[\ell]$. Put $\mathcal S:=\operatorname{span}_\bbQ(f_0,f_1,f_2,y_3)\subset\bbK[\ell]$. An element of $\mathcal S$ lies in $\bbK$ if and only if its coefficient of $y_3$ vanishes, so $\mathcal S\cap\bbK=\operatorname{span}(f_0,f_1,f_2)$, the space of power-series solutions of $\PFfour$.
\item For $\cP=\sum_{k\leq4}a_k\Dq^k$ with coefficients $a_k$ in a differential ring $(\cR_0,\Dq)$, for instance $\cP\in\bbK[\Dq]$, the \emph{formal adjoint} is $\cP^*:=\sum_k(-\Dq)^k\circ a_k$, and the \emph{concomitant} of $\cP$ is the bilinear differential expression
\[
\Conc_\cP(f,g):=\sum_{k=1}^4\sum_{j=0}^{k-1}(-1)^j\,\Dq^j(a_kg)\,\Dq^{k-1-j}f .
\]
We write $\Conc:=\Conc_{\PFfour}$ and expand it as $\Conc(f,g)=\sum_{a,b=0}^3\Dq^af\,\Jet_{ab}\,\Dq^bg$ with $\Jet_{ab}\in\bbQ[u]\subset\bbK$; $\Jet=(\Jet_{ab})$ is the \emph{jet matrix}. As the $A_k$ lie in $\bbQ[u]$, $\PFfour$, $\PFfour^*$ and $\Conc$ act on every differential ring containing $(\bbQ[u],\Dq)$: on $\bbK[\ell]$, $\bbQ\fps q[\logq]$ ($u=1-5^5q$), $\Nring$ ($u=L^{-5}$) and $\bbQ\lser s$ (\cref{def:set-confexp,lem:gau-confexp,cor:gau-rowsNN}).
\item For a basis $\mathbf x=(x_0,\dots,x_3)$ of $\mathcal S$, the \emph{Gram matrix} is $\Conc_{\mathbf x}:=(\Conc(x_i,x_j))_{i,j=0}^3$ and the \emph{Wronskian matrix} is $\Wr_{\mathbf x}:=(\Dq^ax_i)_{a,i=0}^3$ (row $a$, column $i$), whose determinant is the Wronskian $\Wr(x_0,\dots,x_3)$ in the sense of \cref{lem:set-wronskian}. The basis $\mathbf f:=(f_0,f_1,f_2,y_3)$ is the standard one.
\item $\sh$ is the $\bbK$-algebra automorphism of $\bbK[\ell]$ with $\sh(\ell)=\ell+1$.
\item For $t\in\bbK$ with $\Dq t\neq0$, $\de_t:=(\Dq t)^{-1}\Dq$, a derivation of $\bbK[\ell]$. For a basis $\mathbf x$ of $\mathcal S$ (with $\Conc_{\mathbf x}$ invertible, \cref{lem:gau-conc}) the \emph{transversality sums} are $\Upsilon_k:=\sum_{i,j}(\Conc_{\mathbf x}^{-1})_{ij}\,x_i\,\Dq^kx_j$ and, for $y_0\in\bbK^\times$, $\widetilde\Upsilon_k:=\sum_{i,j}(\Conc_{\mathbf x}^{-1})_{ij}\,(x_i/y_0)\,\de_t^k(x_j/y_0)$ ($k\geq0$), elements of the fraction field of $\bbK[\ell]$.
\item For $y_0,y_1\in\bbK$ with $y_0\neq0$ and $t:=y_1/y_0$ not constant, the \emph{cyclic operator} is
\begin{equation}
\label{eq:gau-Mcyc}
\Mcyc(y_0,y_1):=\frac1{y_0}\circ\Dq\circ\frac1{\Dq t}\circ\Dq\circ\bigl(u\,y_0^2(\Dq t)^2\bigr)\circ\Dq\circ\frac1{\Dq t}\circ\Dq\circ\frac1{y_0}\ \in\bbK[\Dq].
\end{equation}
\end{enumerate}
\end{defn}

\begin{lem}[The conifold expansion is a ring map]
\label{lem:gau-confexp}
Let $(\alpha,\beta)\in\bbQ^2$, $\varpi_0=f_0+\alpha\varpi_1+\beta f_2$ and $\tc=\varpi_1/\varpi_0$ (\cref{def:set-frames}), and extend $\Dq$ from $\bbK$ to $\bbQ\lser s$, $s^5=u$, by $\Dq s=-(1-s^5)/(5s^4)$. The assignment of \cref{def:set-confexp} on the generators of $\Nring$ extends to a unique ring homomorphism $\Cexp_{\alpha,\beta}\colon\Nring\to\bbQ\lser s$, and $\Cexp_{\alpha,\beta}\circ\Dq=\Dq\circ\Cexp_{\alpha,\beta}$. For the elements $C_0=c_0$, $I_1=c_0\logq+i_0$, $T=I_1/C_0$ and $C_1=\mathsf W/c_0^2$ of $\Nring$ (\cref{def:set-confexp}), which evaluate to the functions of \cref{def:set-genus0},
\[
\Cexp_{\alpha,\beta}(C_0)=\varpi_0,\qquad \Cexp_{\alpha,\beta}(I_1)=\varpi_1,\qquad \Cexp_{\alpha,\beta}(\mathsf W)=\varpi_0\Dq\varpi_1-\varpi_1\Dq\varpi_0=\varpi_0^2\,\Dq\tc ,
\]
so the conifold expansion sends $C_0$, $I_1$, $T$ and $C_1$ to $\varpi_0$, $\varpi_1$, $\tc$ and $\Dq\tc$.
\end{lem}

\begin{thm}[The cyclic form of $\PFfour$, (CYC)]
\label{thm:gau-cyc}
Let $y_0\in\mathcal S\cap\bbK$ be non-zero and not proportional to $\varpi_1$, and let $\varpi_1$ be taken at any non-zero scale. Then, as operators in $\bbK[\Dq]$,
\begin{equation}
\label{eq:gau-cyc}
\PFfour=\Mcyc(y_0,\varpi_1)=\frac1{y_0}\circ\Dq\circ\frac1{\Dq t}\circ\Dq\circ\bigl(u\,y_0^2(\Dq t)^2\bigr)\circ\Dq\circ\frac1{\Dq t}\circ\Dq\circ\frac1{y_0},\qquad t=\frac{\varpi_1}{y_0}.
\end{equation}
More generally $\PFfour=\Mcyc(y_0,y_1)$ for every basis $(y_0,y_1)$ of an isotropic plane in $\mathcal S\cap\bbK$; these are the planes that contain $\varpi_1$ (\cref{lem:gau-iso}~(iii)). $\Mcyc(y_0,y_1)$ is unchanged under $y_0\mapsto\mu y_0$, $y_1\mapsto\mu'y_1$ ($\mu,\mu'\in\bbQ^\times$) and $y_1\mapsto y_1+ay_0$ ($a\in\bbQ$). In particular $\PFfour=\Mcyc(\varpi_0,\varpi_1)$ in every conifold frame of \cref{def:set-frames}.
\end{thm}

\begin{cor}[The cyclic system is $\PFfive$]
\label{cor:gau-cycsys}
Let $(y_0,y_1)$ be as in \cref{thm:gau-cyc}, $t=y_1/y_0$, and
\[
\kappa=(\kappa_0,\dots,\kappa_4):=\Bigl(\Dq t,\ \frac{L^5}{y_0^2(\Dq t)^2},\ \Dq t,\ y_0,\ y_0\Bigr),\qquad L^5:=u^{-1},
\]
units of $\bbK$. Let $(\cR,\Dq)$ be a differential ring extending $(\bbK,\Dq)$ and $z\in\cR$ with $\Dq z=0$.
\begin{enumerate}
\item For $Y_0,\dots,Y_4\in\cR$ the cyclic system
\begin{equation}
\label{eq:gau-system}
z\Dq Y_a=\kappa_aY_{a+1}\qquad(a\in\bbZ/5)
\end{equation}
holds if and only if $y_0Y_0$ solves $\PFfive$ and $Y_{a+1}=z\Dq Y_a/\kappa_a$ for $a=0,\dots,3$.
\item If $z$ is invertible in $\cR$, then for every $w\in\cR$ the same holds with $z\Dq$ replaced by $z\Dq+w$ in \eqref{eq:gau-system} and in $\PFfive$.
\end{enumerate}
\end{cor}

% Section 3, part 2 of 8: the quintic system and the pulled-back model; Theorem U (notation, statement).

%-------------------------------------------------------------------------------
\subsection{The quintic system and the pulled-back model}
\label{sec:gau-systems}

\begin{defn}[Cyclic systems with one simple pole; regular gauges]
\label{def:gau-onepole}
Let $\Bbbk$ be a field of characteristic $0$. A \emph{cyclic system with one simple pole} over $\Bbbk\fps u$ is a system $z\de_uY=AY$ with $A=\sum_{a\in\bbZ/5}b_aE_{a,a+1}$ whose coefficients have the \emph{one-pole shape}: $b_a\in\Bbbk\fps u^\times$ for $a\neq1$ and $ub_1\in\Bbbk\fps u^\times$. Its \emph{residue} is $r:=(ub_1)(0)\in\Bbbk^\times$ and its \emph{product} is $\Bprod:=\prod_ab_a\in\Bbbk\lser u$. It is \emph{self-dual} if $b_0=b_2$ and $b_3=b_4$, i.e.\ $\eta A=A^T\eta$ for $\eta=(5\,\delta_{a+b\equiv3})_{a,b}$ (the entries $(c,d)$ are $5b_{d-1}\delta_{c+d\equiv4}$ and $5b_{c-1}\delta_{c+d\equiv4}$). For a second such system $\tilde A=\sum_a\tilde b_aE_{a,a+1}$, with residue $\tilde r$, the condition (EP) of \emph{equal products} is $\prod_ab_a=\prod_a\tilde b_a$. The \emph{gauge equation} from $\tilde A$ to $A$ is $zP'=AP-P\tilde A$ for $P=\sum_{k\geq0}z^kP_k$, $P_k\in\mathfrak{gl}_5(\Bbbk\fps u)$; then $z\de_u(P\tilde Y)=AP\tilde Y$ whenever $z\de_u\tilde Y=\tilde A\tilde Y$. $\Gau(A,\tilde A)$ denotes the set of all solutions $P\in\mathfrak{gl}_5(\Bbbk\fps u)\fps z$; it is a $\Bbbk\fps z$-module, since $z(P\cn)'=(AP-P\tilde A)\cn$ for $\cn\in\Bbbk\fps z$. A \emph{regular gauge} is an element of $\Gau(A,\tilde A)$ with $P_0$ invertible. $P$ is \emph{graded} if $P_k=\sum_bp_{k,b}E_{b+k,b}$ for every $k$.
\end{defn}

\begin{prop}[The quintic cyclic system]
\label{prop:gau-quintic}
Let $(\alpha,\beta)\in\bbQ^2$ give a conifold frame (\cref{def:set-frames}), $\tc=\varpi_1/\varpi_0$, and
\[
\kappa:=\Bigl(\Dq\tc,\ \frac{L^5}{\varpi_0^2(\Dq\tc)^2},\ \Dq\tc,\ \varpi_0,\ \varpi_0\Bigr),\qquad L^5=u^{-1},
\]
the vector of \cref{cor:gau-cycsys} for $(y_0,y_1)=(\varpi_0,\varpi_1)$.
\begin{enumerate}
\item $\kappa_a=\Cexp_{\alpha,\beta}(C_1\kappa^{\mathrm{LR}}_a)$ for $a\in\bbZ/5$, where $C_1\kappa^{\mathrm{LR}}=(C_1,L^5/(C_0^2C_1^2),C_1,C_0,C_0)$ (\cref{prop:set-cyclicLR}) is represented in $\Nring$ by $(\mathsf W/c_0^2,\ L^5c_0^2/\mathsf W^2,\ \mathsf W/c_0^2,\ c_0,\ c_0)$. Hence the conifold expansion carries the forward relations $(z\Dq+\zeta^iL)X_a=C_1\kappa^{\mathrm{LR}}_aX_{a+1}$ ($a=0,\dots,3$) of elements of $\Nring\lser z\otimes\bbQ(\zeta)$ to the relations $(z\Dq+\zeta^iL)\Cexp(X_a)=\kappa_a\Cexp(X_{a+1})$, $L=s^{-1}$; the closing relation is then given by \cref{cor:gau-cycsys}.
\item The cyclic system \eqref{eq:gau-system} is $z\de_uY=AY$ with $A=\sum_ab_aE_{a,a+1}$, $b_a:=-\kappa_a/(1-u)$,
\[
b=\Bigl(\tc',\ -\frac1{u\varpi_0^2(1-u)^3\tc'^2},\ \tc',\ -\frac{\varpi_0}{1-u},\ -\frac{\varpi_0}{1-u}\Bigr),
\]
a self-dual cyclic system with one simple pole over $\bbQ\fps u$, with residue $r=-1$, product $\prod_ab_a=-1/(u(1-u)^5)$, and $b_a(0)=1,1,-1,-1$ for $a=0,2,3,4$.
\item More generally, for $y_0\in\mathcal S\cap\bbK$ not proportional to $\varpi_1$ and $t=\varpi_1/y_0$, the coefficients $-\kappa_a/(1-u)$ of \cref{cor:gau-cycsys} have the one-pole shape if and only if $y_0(0)\neq0$.
\end{enumerate}
\end{prop}

\begin{proof}
(i) By \cref{lem:gau-confexp}, $\Cexp(c_0)=\varpi_0$, $\Cexp(\mathsf W/c_0^2)=\Dq\tc$ and $\Cexp(L^5)=u^{-1}$, which gives the five entries (for $\kappa_1$, $u^{-1}\varpi_0^2/(\varpi_0^4(\Dq\tc)^2)$). As $\Cexp$ is a ring homomorphism commuting with $\Dq$, with $\Cexp(L)=s^{-1}$, it carries each forward relation to the relation with $\kappa_a$; the closing relation is \cref{cor:gau-cycsys}~(ii) with $w=\zeta^is^{-1}$. (ii) As $\Dq=-(1-u)\,\dd/\dd u$, $z\Dq Y_a=\kappa_aY_{a+1}$ is $z\de_uY_a=b_aY_{a+1}$, and $\Dq\tc=-(1-u)\tc'$, $L^5=u^{-1}$ give $b$. As $\varpi_0(0)=\tc'(0)=1$ (\cref{def:set-frames}), $b_0,b_2,b_3,b_4$ are units with the stated values and $(ub_1)(0)=-1$; self-duality is visible, and $\prod_ab_a=\tc'^2\cdot\bigl(-u^{-1}\varpi_0^{-2}(1-u)^{-3}\tc'^{-2}\bigr)\cdot\varpi_0^2(1-u)^{-2}$. (iii) If $y_0(0)\neq0$, then $t=u/y_0(0)+O(u^2)$, so $t'(0)\neq0$ and the computation of (ii) applies to $(y_0,t)$; if $y_0(0)=0$, $-y_0/(1-u)$ is not a unit.
\end{proof}

\begin{prop}[Equal products and $v(u)$]
\label{prop:gau-equalprod}
Let $v\in u\bbQ\fps u$ with $v'(0)\neq0$. The model system $z\,\dd Y_b/\dd v=b^{\mathrm m}_bY_{b+1}$, $b^{\mathrm m}=(1,-v^{-1},1,-1,-1)$ (\cref{prop:set-modrows}), pulled back along $v=v(u)$ is $z\de_uY=\tilde AY$ with $\tilde A=v'(u)\sum_bb^{\mathrm m}_b(v(u))E_{b,b+1}$, i.e.
\[
\tilde b=\bigl(v',\ -v'/v,\ v',\ -v',\ -v'\bigr),\qquad \prod_a\tilde b_a=-v'^5/v .
\]
\begin{enumerate}
\item If $v=u+O(u^2)$, $\tilde A$ is a self-dual cyclic system with one simple pole over $\bbQ\fps u$, with residue $\tilde r=-1$ and $\tilde b_a(0)=1,1,-1,-1$ for $a=0,2,3,4$.
\item For $v=u+O(u^2)$, the products of $\tilde A$ and of the quintic system of \cref{prop:gau-quintic} are equal if and only if $v'^5/v=1/(u(1-u)^5)$, if and only if $v=u\,\tau(u)^{-5/4}$ (\cref{def:set-tau}).
\item (Necessity of (EP).) If $A,\tilde A$ are cyclic systems with one simple pole over $\Bbbk\fps u$ and $\Gau(A,\tilde A)\neq0$, then $\prod_ab_a=\prod_a\tilde b_a$.
\end{enumerate}
In particular, for $v=v(u)$ and every conifold frame, the quintic system and the pulled-back model are two self-dual cyclic systems with one simple pole satisfying (EP), and $b_a(0)=\tilde b_a(0)$ for $a\neq1$, $r=\tilde r=-1$.
\end{prop}

\begin{proof}
The chain rule gives $\tilde A$, and $\prod_a\tilde b_a=v'^5\cdot(-v^{-1})\cdot(-1)^2$. (i) $v'$ is a unit with $v'(0)=1$, and $-uv'/v$ a unit with value $-1$ at $u=0$. (ii) The first equivalence is \cref{prop:gau-quintic}~(ii). Write $v=uw$, $w\in1+u\bbQ\fps u$; elements of $1+u\bbQ\fps u$ have unique $n$-th roots there (binomial series; $c^n\neq1$ for $c\in1+u\bbQ\fps u\smallsetminus\{1\}$). As $v'(0)=1$, $v'^5=w(1-u)^{-5}$ is equivalent to $v'=w^{1/5}(1-u)^{-1}$. In $\bbQ\lser s$, $s^5=u$, with $v^{4/5}:=s^4w^{4/5}$ and $\frac{\dd}{\dd u}v^{4/5}=\frac45s^{-1}w^{-1/5}v'$, this reads $\frac{\dd}{\dd u}v^{4/5}=\frac45s^{-1}(1-u)^{-1}$, which is also the derivative of $s^4/\tau=\sum_{n\geq0}\frac4{5n+4}u^{n+4/5}$ (\cref{def:set-tau}). As $\dd/\dd u$ has kernel $\bbQ$ on $\bbQ\lser s$ and $v^{4/5}-s^4/\tau\in s^4\bbQ\fps u$, the equation is equivalent to $w^{4/5}=\tau^{-1}$, i.e.\ $w=\tau^{-5/4}$. (iii) If $P\neq0$ and $P_n$ is its first non-zero coefficient, the coefficient of $z^n$ of the gauge equation is $AP_n=P_n\tilde A$, so $A^5P_n=P_n\tilde A^5$; as $A^5=(\prod_ab_a)1$ and $\tilde A^5=(\prod_a\tilde b_a)1$, (EP) follows. The last sentence combines (i), (ii) and \cref{prop:gau-quintic}~(ii).
\end{proof}

%-------------------------------------------------------------------------------
\subsection{Theorem U}
\label{sec:gau-U}

Throughout this subsection $\Bbbk$ is a field of characteristic $0$, $A=\sum_ab_aE_{a,a+1}$ and $\tilde A=\sum_a\tilde b_aE_{a,a+1}$ are cyclic systems with one simple pole over $\Bbbk\fps u$ satisfying (EP), and a prime is $\dd/\dd u$. The following notation is local to this subsection.
\begin{itemize}
\item $\mathfrak{gl}^{(m)}:=\bigoplus_a\Bbbk\lser uE_{a,a+m}$ for $m\in\bbZ/5$; an element of $\mathfrak{gl}^{(m)}$ has \emph{shift} $m$, and $\mathfrak{gl}_5(\Bbbk\lser u)=\bigoplus_m\mathfrak{gl}^{(m)}$.
\item $\Bprod^{(m)}_a:=\tilde b_a\tilde b_{a+1}\cdots\tilde b_{a+m-1}$ for $0\leq m\leq5$ (so $\Bprod^{(0)}_a=1$ and $\Bprod^{(5)}_a=\Bprod:=\prod_a\tilde b_a$), and $\mathcal Y:=\tilde A^{-1}$.
\item $\Theta_U:=\diag(d^U_0,\dots,d^U_4)$ with $d^U_0:=1$ and $d^U_{a+1}:=d^U_a\tilde b_a/b_a$; $\operatorname{dlog}\Theta_U:=\Theta_U^{-1}\Theta_U'$; $\mathfrak d:=(\det\Theta_U/\det\Theta_U(0))^{-1/5}$.
\item If $\Bbbk$ contains $\zeta$, put $\mathrm{gr}:=\diag(0,1,2,3,4)$, $\zeta^{i(\mathrm{gr}+1)}:=\diag(\zeta^{i(a+1)})_a$, and define the \emph{colour rotation} on matrices and on vectors by
\[
\operatorname{Rot}_iP(u,z):=\zeta^{i(\mathrm{gr}+1)}P(u,\zeta^{-i}z)\zeta^{-i(\mathrm{gr}+1)},\qquad \operatorname{Rot}_iY(u,z):=\zeta^{i(\mathrm{gr}+1)}Y(u,\zeta^{-i}z).
\]
\end{itemize}
By (EP), $d^U_4\tilde b_4/b_4=\prod_a\tilde b_a/b_a=1=d^U_0$, so the recursion of $d^U$ holds for every $a\in\bbZ/5$. The $d^U_a$ are units of $\Bbbk\fps u$ ($\tilde b_1/b_1=(u\tilde b_1)/(ub_1)$), so $\mathfrak d\in1+u\Bbbk\fps u$ is a binomial series and $\mathfrak d'=-\frac15\tr(\operatorname{dlog}\Theta_U)\,\mathfrak d$; and $(A\Theta_U)_{a,a+1}=b_ad^U_{a+1}=d^U_a\tilde b_a=(\Theta_U\tilde A)_{a,a+1}$, so $A\Theta_U=\Theta_U\tilde A$.

\begin{thm}[Theorem U]
\label{thm:gau-U}
Let $A$, $\tilde A$ be cyclic systems with one simple pole over $\Bbbk\fps u$ with equal products.
\begin{enumerate}
\item[(a)] (Existence.) There is $P^\natural\in\Gau(A,\tilde A)$ with $P^\natural_0=\mathfrak d\,\Theta_U$, a diagonal matrix with unit entries.
\item[(b)] (Uniqueness.) For every $P^\natural\in\Gau(A,\tilde A)$ with $P^\natural_0=\mathfrak d\,\Theta_U$, the map $\Bbbk\fps z\to\Gau(A,\tilde A)$, $\cn\mapsto P^\natural\cn(z)$, is bijective. Consequently every solution $P$ of the gauge equation with coefficients in $\mathfrak{gl}_5(\Bbbk\fps u)$ is $P^\natural\cn(z)$ for a unique $\cn\in\Bbbk\fps z$; its $z^0$-term is $P_0=\cn(0)\,\mathfrak d\,\Theta_U$, a diagonal matrix, and $P$ is a regular gauge if and only if $\cn(0)\neq0$. If $P$ is one regular gauge, the regular gauges are exactly the $P\cn(z)$ with $\cn\in\Bbbk\fps z^\times$. For the pair $(\tilde A,\tilde A)$: $\Gau(\tilde A,\tilde A)=\Bbbk\fps z\cdot1$.
\item[(c)] (Graded gauges.) Graded regular gauges exist; if $P$ is one, the graded regular gauges are exactly the $P\cn(z)$ with $\cn\in\Bbbk\fps{z^5}^\times$.
\item[(d)] (The field.) If all $b_a$, $\tilde b_a$ lie in $\bbQ\lser u$, then $P^\natural$ can be chosen in $\mathfrak{gl}_5(\bbQ\fps u)\fps z$; for every field $\Bbbk\supseteq\bbQ$ the solutions with coefficients in $\mathfrak{gl}_5(\Bbbk\fps u)$ are then $P^\natural\Bbbk\fps z$, and $P^\natural\cn$ has rational coefficients if and only if $\cn\in\bbQ\fps z$.
\end{enumerate}
\end{thm}

% Section 3, part 3 of 8: the commutant, lifting, the scalar equation.

\begin{lem}[The commutant]
\label{lem:gau-comm}
Let $F$ be a field of characteristic $0$, $\tilde b_0,\dots,\tilde b_4\in F^\times$, $\tilde A=\sum_a\tilde b_aE_{a,a+1}$, and define $\mathfrak{gl}^{(m)}$, $\Bprod^{(m)}_a$, $\Bprod$ and $\mathcal Y$ over $F$ as above.
\begin{enumerate}
\item[(a)] $\tilde A^m=\sum_a\Bprod^{(m)}_aE_{a,a+m}$ for $0\leq m\leq5$, so $\tilde A^5=\Bprod\,1$; $\mathcal Y=\sum_a\tilde b_a^{-1}E_{a+1,a}$, $\mathcal Y^j=\sum_a(\Bprod^{(j)}_a)^{-1}E_{a+j,a}$ for $0\leq j\leq5$, and $\tilde A^m=\Bprod\,\mathcal Y^{5-m}$ for $1\leq m\leq5$; $\det(x1-\tilde A)=x^5-\Bprod$.
\item[(b)] For $X=\sum_a\xi_aE_{a,a+m}\in\mathfrak{gl}^{(m)}$, $[\tilde A,X]=\sum_a(\tilde b_a\xi_{a+1}-\xi_a\tilde b_{a+m})E_{a,a+m+1}\in\mathfrak{gl}^{(m+1)}$.
\item[(c)] For every $m$, $\ker(\operatorname{ad}\tilde A)\cap\mathfrak{gl}^{(m)}=F\tilde A^m$; hence $\ker(\operatorname{ad}\tilde A)=F[\tilde A]=F[\mathcal Y]=\bigoplus_{j=0}^4F\mathcal Y^j$, of dimension $5$.
\item[(d)] For $m\in\{1,\dots,5\}$ and $X=\sum_a\xi_aE_{a,a+m}\in\mathfrak{gl}^{(m)}$ put $\lambda_m(X):=\sum_a\xi_a/\Bprod^{(m)}_a$ (a local letter). Then $\lambda_m(\tilde A^m)=5$, $\operatorname{ad}\tilde A(\mathfrak{gl}^{(m-1)})=\ker\lambda_m$, $\mathfrak{gl}^{(m)}=F\tilde A^m\oplus\ker\lambda_m$, and $\mathfrak{gl}_5(F)=F[\mathcal Y]\oplus\operatorname{im}(\operatorname{ad}\tilde A)$.
\item[(e)] The projection $\operatorname{pr}$ onto $F[\mathcal Y]$ along $\operatorname{im}(\operatorname{ad}\tilde A)$ is $F$-linear and preserves the shift; for $X\in\mathfrak{gl}^{(m)}$, $1\leq m\leq5$,
\[
\operatorname{pr}(X)=\tfrac15\lambda_m(X)\,\tilde A^m=\tfrac15\Bigl[\sum_a\xi_a\Bprod^{(5-m)}_{a+m}\Bigr]\mathcal Y^{5-m}.
\]
Equivalently, $\operatorname{pr}(X)$ is the unique element of $F[\mathcal Y]$ with $\tr(\operatorname{pr}(X)Z)=\tr(XZ)$ for all $Z\in F[\mathcal Y]$.
\item[(f)] Let $F=\Bbbk\lser u$ and let $\tilde A$ have the one-pole shape. Then $\mathcal Y\in\mathfrak{gl}_5(\Bbbk\fps u)$; an element $\sum_{j=0}^4w_j\mathcal Y^j$ ($w_j\in\Bbbk\lser u$) lies in $\mathfrak{gl}_5(\Bbbk\fps u)$ if and only if every $w_j$ lies in $\Bbbk\fps u$; the commutant of $\tilde A$ in $\mathfrak{gl}_5(\Bbbk\fps u)$ is the ring $\Bbbk\fps u[\mathcal Y]:=\bigoplus_{j=0}^4\Bbbk\fps u\,\mathcal Y^j$; and for $X\in\mathfrak{gl}_5(\Bbbk\fps u)$ the coefficient of $\mathcal Y^j$ in $\operatorname{pr}(X)$ lies in $u^{-1}\Bbbk\fps u$ for $j=1,\dots,4$ and in $\Bbbk\fps u$ for $j=0$.
\item[(g)] (Not used.) $x^5-\Bprod$ is separable, so $\tilde A$ is regular semisimple; in the situation of (f), $x^5-\Bprod$ is irreducible over $\Bbbk\lser u$, so $F[\mathcal Y]$ is a field.
\end{enumerate}
\end{lem}

\begin{proof}
(a) Induction on $m$, with $\Bprod^{(m)}_a\tilde b_{a+m}=\Bprod^{(m+1)}_a$ and $E_{a,a+5}=E_{aa}$; $\tilde A\sum_c\tilde b_c^{-1}E_{c+1,c}=1$ gives $\mathcal Y$, and $\mathcal Y^j$ follows likewise. Reindexing, $\mathcal Y^{5-m}=\sum_a(\Bprod^{(5-m)}_{a+m})^{-1}E_{a,a+m}$, and $\Bprod^{(m)}_a\Bprod^{(5-m)}_{a+m}=\Bprod$. In the Leibniz expansion of $\det(x1-\tilde A)$ only the identity and the even $5$-cycle $a\mapsto a+1$ contribute ($\pi(a)=a+1$ forces $\pi(a+1)=a+2$, and so on), the latter with $\prod_a(-\tilde b_a)=-\Bprod$.
(b) $\tilde AX=\sum_a\tilde b_a\xi_{a+1}E_{a,a+1+m}$ and $X\tilde A=\sum_a\xi_a\tilde b_{a+m}E_{a,a+m+1}$.
(c) By (b) $\operatorname{ad}\tilde A$ shifts the (direct) grading by $1$, so it suffices to treat $X\in\mathfrak{gl}^{(m)}$: $[\tilde A,X]=0$ means $\xi_{a+1}=\xi_a\tilde b_{a+m}/\tilde b_a$, whose solutions are determined by $\xi_0$, and $\xi_a=\Bprod^{(m)}_a$ is one. With $\tilde A^m=\Bprod\mathcal Y^{5-m}$ and $\mathcal Y=\Bprod^{-1}\tilde A^4$, $F[\tilde A]=F[\mathcal Y]$.
(d) For $Z\in\ker(\operatorname{ad}\tilde A)$, $\tr(Z[\tilde A,Y])=0$. For $X\in\mathfrak{gl}^{(m)}$, (a) gives $\tr(X\tilde A^{5-m})=\sum_a\xi_a\Bprod^{(5-m)}_{a+m}=\Bprod\lambda_m(X)$, so $\lambda_m$ vanishes on $\operatorname{ad}\tilde A(\mathfrak{gl}^{(m-1)})$, which has dimension $4$ by (c); as $\lambda_m(\tilde A^m)=5\neq0$, this image is $\ker\lambda_m$. Summing over $m$ gives the last assertion.
(e) $X=\frac15\lambda_m(X)\tilde A^m+\bigl(X-\frac15\lambda_m(X)\tilde A^m\bigr)$ with the second term in $\ker\lambda_m$; the second form uses (a). $X-\operatorname{pr}(X)\in\operatorname{im}(\operatorname{ad}\tilde A)$ is trace-orthogonal to $F[\mathcal Y]$ by (d), and the trace form is non-degenerate on $F[\tilde A]$: in the basis $1,\tilde A,\dots,\tilde A^4$ its Gram matrix has $5$ at $(0,0)$, $5\Bprod$ at $(i,5-i)$ and $0$ elsewhere ($\tilde A^n$ has shift $n$).
(f) The entries of $\mathcal Y$ are the $\tilde b_a^{-1}$, integral since $\tilde b_1^{-1}=u(u\tilde b_1)^{-1}$. The supports $\{(a+j,a)\}$ of $\mathcal Y^0,\dots,\mathcal Y^4$ partition the $25$ positions, and the entry $(a+j,a)$ of $\sum_jw_j\mathcal Y^j$ is $w_j(\Bprod^{(j)}_a)^{-1}$. For $j\leq4$ the window $\{a,\dots,a+j-1\}$ contains the pole index $1$ at most once, so $(\Bprod^{(j)}_a)^{-1}$ is integral, and a unit for $a=2$; this gives the criterion. $\Bbbk\fps u[\mathcal Y]$ is a ring since $\mathcal Y^5=\Bprod^{-1}1$, $\Bprod^{-1}\in u\Bbbk\fps u$, and by (c) it is the integral commutant. For integral $X\in\mathfrak{gl}^{(m)}$, $1\leq m\leq4$, the coefficient $\frac15\sum_a\xi_a\Bprod^{(5-m)}_{a+m}$ involves at most four consecutive weights, at most one of them $\tilde b_1$, so it lies in $u^{-1}\Bbbk\fps u$; for $m=5$ it is $\frac15\tr X$.
(g) $\gcd(x^5-\Bprod,5x^4)=1$. In the situation of (f), with $\Bprod^\circ:=u\Bprod$ (a unit) and $x=1/X$, $x^5-\Bprod=-\Bprod^\circ u^{-1}X^{-5}\bigl(X^5-u(\Bprod^\circ)^{-1}\bigr)$, an Eisenstein polynomial at $u$.
\end{proof}

\begin{lem}[Lifting]
\label{lem:gau-lift}
Let $\tilde A$ have the one-pole shape, let $X\in\mathfrak{gl}_5(\Bbbk\fps u)$ have shift $m+1$ with $0\leq m\leq4$, and suppose $\operatorname{pr}(X)=0$. Then there is $\mathcal Z\in\mathfrak{gl}_5(\Bbbk\fps u)$ of shift $m$ with $[\tilde A,\mathcal Z]=X$, and every solution $\mathcal Z'\in\mathfrak{gl}_5(\Bbbk\fps u)$ of $[\tilde A,\mathcal Z']=X$ is $\mathcal Z+w$ with $w\in\Bbbk\fps u[\mathcal Y]$. If $\tilde A$ and $X$ have rational coefficients, $\mathcal Z$ can be chosen rational.
\end{lem}

\begin{proof}
Write $X=\sum_a\xi_aE_{a,a+m+1}$ and try $\mathcal Z=\sum_a\rho_a\Bprod^{(m)}_aE_{a,a+m}$ (local letters). By \cref{lem:gau-comm}~(b), $[\tilde A,\mathcal Z]=\sum_a(\rho_{a+1}-\rho_a)\Bprod^{(m+1)}_aE_{a,a+m+1}$, so the equation is $\rho_{a+1}-\rho_a=\sigma_a:=\xi_a/\Bprod^{(m+1)}_a$, solvable if and only if $\sum_a\sigma_a=\lambda_{m+1}(X)=0$, i.e.\ $\operatorname{pr}(X)=0$ (\cref{lem:gau-comm}~(e)). Each $\sigma_a$ is integral, with $\sigma_a(0)=0$ when the window $\{a,\dots,a+m\}$ contains $1$. Put $\rho_1:=0$, $\rho_{a+1}:=\rho_a+\sigma_a$ ($a=1,\dots,4$); the relation for $a=0$ follows from $\sum_a\sigma_a=0$. The entry $\rho_a\Bprod^{(m)}_a$ is integral when the window $\{a,\dots,a+m-1\}$ avoids $1$, and otherwise ($a=2-m,\dots,1$) when $\rho_a(0)=0$: for $a=1$ by choice, and for $2-m\leq a\leq0$ because $\rho_a=-\sum_{c=a}^0\sigma_c$ with every window $\{c,\dots,c+m\}$ containing $1$. Two integral solutions differ by an element of $\Bbbk\fps u[\mathcal Y]$ (\cref{lem:gau-comm}~(c), (f)); the construction uses only field operations.
\end{proof}

\begin{lem}[The scalar equation]
\label{lem:gau-scal}
Let $A$, $\tilde A$ have the one-pole shape and satisfy (EP).
\begin{enumerate}
\item[(a)] For $w=\sum_{j=0}^4w_j\mathcal Y^j$ with $w_j\in\Bbbk\lser u$,
\[
\operatorname{pr}\bigl(w'+\operatorname{dlog}\Theta_U\,w\bigr)=\sum_{j=0}^4\bigl(w_j'+\chi^U_jw_j\bigr)\mathcal Y^j,\qquad \chi^U_j:=-\frac j5\,\frac{\Bprod'}{\Bprod}+\frac15\tr(\operatorname{dlog}\Theta_U).
\]
In particular $\operatorname{pr}((\mathcal Y^j)')=-\frac j5(\Bprod'/\Bprod)\mathcal Y^j$ and $\operatorname{pr}(\operatorname{dlog}\Theta_U\,\mathcal Y^j)=\frac15\tr(\operatorname{dlog}\Theta_U)\,\mathcal Y^j$. Moreover $\chi^U_j=\frac j{5u}+\chi^\circ_j$ with $\chi^\circ_j:=-\frac j5(\Bprod^\circ)'/\Bprod^\circ+\frac15\tr(\operatorname{dlog}\Theta_U)\in\Bbbk\fps u$, where $\Bprod^\circ:=u\Bprod\in\Bbbk\fps u^\times$.
\item[(b)] For $Z\in\mathfrak{gl}_5(\Bbbk\fps u)$ the equation $\operatorname{pr}(w'+\operatorname{dlog}\Theta_U\,w)=-\operatorname{pr}(Z)$ has solutions $w\in\Bbbk\fps u[\mathcal Y]$; their components $w_1,\dots,w_4$ are unique, and $w_0$ is unique up to $w_0\mapsto w_0+\cn\,\mathfrak d$, $\cn\in\Bbbk$.
\item[(c)] For $Z=0$ the solutions in $\Bbbk\fps u[\mathcal Y]$ are exactly the $\cn\,\mathfrak d\,1$, $\cn\in\Bbbk$.
\item[(d)] If the data are rational, so are the solutions for $\cn\in\bbQ$; $\mathfrak d\in1+u\bbQ\fps u$.
\end{enumerate}
\end{lem}

\begin{proof}
(a) $(\mathcal Y^j)'$ has shift $5-j$, with entries $((\Bprod^{(j)}_a)^{-1})'$ at $(a+j,a)$; by \cref{lem:gau-comm}~(e),
\[
\operatorname{pr}((\mathcal Y^j)')=\frac15\Bigl[\sum_a\bigl((\Bprod^{(j)}_a)^{-1}\bigr)'\Bprod^{(j)}_a\Bigr]\mathcal Y^j=-\frac15\Bigl[\sum_a\frac{(\Bprod^{(j)}_a)'}{\Bprod^{(j)}_a}\Bigr]\mathcal Y^j ,
\]
and $\sum_a(\Bprod^{(j)}_a)'/\Bprod^{(j)}_a=\sum_a\sum_{i<j}\tilde b_{a+i}'/\tilde b_{a+i}=j\Bprod'/\Bprod$, each index occurring $j$ times. Likewise $\operatorname{dlog}\Theta_U\mathcal Y^j$ has entries $\frac{(d^U_{a+j})'}{d^U_{a+j}}(\Bprod^{(j)}_a)^{-1}$, and \cref{lem:gau-comm}~(e) gives $\frac15\tr(\operatorname{dlog}\Theta_U)\mathcal Y^j$. As $\operatorname{pr}$ is $\Bbbk\lser u$-linear and fixes $\Bbbk\lser u[\mathcal Y]$, the formula follows. Finally $\Bprod'/\Bprod=(\Bprod^\circ)'/\Bprod^\circ-u^{-1}$, and $\tr(\operatorname{dlog}\Theta_U)$ is integral since the $d^U_c$ are units.

(b) By \cref{lem:gau-comm}~(f), $\operatorname{pr}(Z)=\sum_j\phi^Z_j\mathcal Y^j$ with $u\phi^Z_j\in\Bbbk\fps u$ for $j\geq1$ and $\phi^Z_0\in\Bbbk\fps u$ (local letters), and by (a) the equation splits into $w_j'+\chi^U_jw_j=-\phi^Z_j$. For $j=1,\dots,4$, multiplied by $u$, the coefficient of $u^n$ reads
\[
\Bigl(n+\frac j5\Bigr)w_{j,n}=-\coef{u^n}(u\phi^Z_j)-\sum_{l=0}^{n-1}\coef{u^{n-1-l}}\chi^\circ_j\cdot w_{j,l},
\]
with $n+j/5\neq0$: exactly one solution $w_j\in\Bbbk\fps u$. For $j=0$, $\chi^U_0=-\mathfrak d'/\mathfrak d$, so the equation is $\mathfrak d\,(w_0/\mathfrak d)'=-\phi^Z_0$, with solutions $w_0=\mathfrak d\bigl(\cn-\int_0^u\phi^Z_0/\mathfrak d\,\dd u\bigr)$, $\cn\in\Bbbk$ ($\int_0^u$ the formal antiderivative without constant term), and no others. (c) is (b) with $\phi^Z=0$. (d) The pivots are rational and $\mathfrak d$ is a binomial series of an element of $1+u\bbQ\fps u$.
\end{proof}

% Section 3, part 4 of 8: proof of Theorem U, the leading term, the Z/5-degree, graded gauges.

\begin{proof}[Proof of \cref{thm:gau-U}~(a)]
Write $P=\Theta_U\check P$ (a local letter); $\check P$ is integral if and only if $P$ is. As $A\Theta_U=\Theta_U\tilde A$, the gauge equation is equivalent to $z\check P'=[\tilde A,\check P]-z\operatorname{dlog}\Theta_U\,\check P$, i.e.\ to
\begin{equation}
\label{eq:gau-reduced}
[\tilde A,\check P_0]=0,\qquad [\tilde A,\check P_{k+1}]=\check P_k'+\operatorname{dlog}\Theta_U\,\check P_k\quad(k\geq0).
\end{equation}
We construct $\check P_k=\mathcal Z_k+w^{(k)}$ with $\mathcal Z_k$ integral and $w^{(k)}\in\Bbbk\fps u[\mathcal Y]$ such that (I$_k$) $[\tilde A,\check P_k]=\check P_{k-1}'+\operatorname{dlog}\Theta_U\check P_{k-1}$ ($\check P_{-1}:=0$) and (S$_k$) $\operatorname{pr}(\check P_k'+\operatorname{dlog}\Theta_U\check P_k)=0$. For $k=0$ take $\mathcal Z_0:=0$, $w^{(0)}:=\mathfrak d\,1$: (I$_0$) is clear and (S$_0$) is \cref{lem:gau-scal}~(c). Given $\check P_k$, $X_{k+1}:=\check P_k'+\operatorname{dlog}\Theta_U\check P_k$ is integral with $\operatorname{pr}(X_{k+1})=0$ by (S$_k$), so \cref{lem:gau-lift}, applied to each shift component, gives an integral $\mathcal Z_{k+1}$ with $[\tilde A,\mathcal Z_{k+1}]=X_{k+1}$; then \cref{lem:gau-scal}~(b) with $Z:=\mathcal Z_{k+1}'+\operatorname{dlog}\Theta_U\mathcal Z_{k+1}$ gives $w^{(k+1)}$, and (I$_{k+1}$) holds since $[\tilde A,w^{(k+1)}]=0$, (S$_{k+1}$) since its left side is $\operatorname{pr}(Z)+\operatorname{pr}\bigl((w^{(k+1)})'+\operatorname{dlog}\Theta_Uw^{(k+1)}\bigr)=0$. By \eqref{eq:gau-reduced}, $P^\natural:=\Theta_U\sum_kz^k\check P_k\in\Gau(A,\tilde A)$ and $P^\natural_0=\mathfrak d\,\Theta_U$.
\end{proof}

\begin{lem}[The leading term of a solution of the gauge equation]
\label{lem:gau-lead}
Let $A$, $\tilde A$ have the one-pole shape and satisfy (EP), let $P\in\Gau(A,\tilde A)$ and $n\geq0$ with $P_0=\dots=P_{n-1}=0$. Then $P_n=\cn\,\mathfrak d\,\Theta_U$ for a unique $\cn\in\Bbbk$, namely $\cn=(P_n)_{00}(0)$.
\end{lem}

\begin{proof}
For $\check P:=\Theta_U^{-1}P$, the coefficients of $z^n$ and $z^{n+1}$ of \eqref{eq:gau-reduced} give $[\tilde A,\check P_n]=0$, so $\check P_n\in\Bbbk\fps u[\mathcal Y]$ (\cref{lem:gau-comm}~(c), (f)), and $\check P_n'+\operatorname{dlog}\Theta_U\check P_n\in\operatorname{im}(\operatorname{ad}\tilde A)$, so its projection vanishes. By \cref{lem:gau-scal}~(c), $\check P_n=\cn\,\mathfrak d\,1$, and $\cn=(P_n)_{00}(0)$ as $d^U_0=\mathfrak d(0)=1$.
\end{proof}

\begin{proof}[Proof of \cref{thm:gau-U}~(b)--(d)]
Part (c) is \cref{prop:gau-graded}. (b) \emph{Injectivity:} if $\cn\neq0$ has order $n$, $\coef{z^n}(P^\natural\cn)=\cn_n\mathfrak d\Theta_U\neq0$. \emph{Surjectivity:} given $P\in\Gau(A,\tilde A)$, suppose $\cn_0,\dots,\cn_{n-1}$ are found with $R^{(n)}:=P-P^\natural\sum_{m<n}\cn_mz^m\equiv0$ modulo $z^n$. As $\Gau(A,\tilde A)$ is a $\Bbbk\fps z$-module, \cref{lem:gau-lead} gives $\coef{z^n}R^{(n)}=\cn_n\mathfrak d\Theta_U$ for a unique $\cn_n\in\Bbbk$, and then $R^{(n+1)}=R^{(n)}-\cn_nz^nP^\natural\equiv0$ modulo $z^{n+1}$. With $\cn:=\sum_n\cn_nz^n$, every coefficient of $P-P^\natural\cn$ vanishes. Its $z^0$-term is $\cn(0)\mathfrak d\Theta_U$; the statements on regular gauges follow, since $P^\natural\cn'=P\cdot(\cn_P^{-1}\cn')$ for $P=P^\natural\cn_P$ regular. For $(\tilde A,\tilde A)$, $\Theta_U=\mathfrak d=1$ and $1\in\Gau(\tilde A,\tilde A)$.

(d) Every object in the proof of (a) is obtained from the $b_a$, $\tilde b_a$ by field operations, derivatives and $\int_0^u$; with rational free constants the construction stays in $\mathfrak{gl}_5(\bbQ\fps u)\fps z$. Over $\Bbbk\supseteq\bbQ$ the proofs of \cref{lem:gau-comm,lem:gau-lift,lem:gau-scal,lem:gau-lead} hold verbatim, so (b) applies; and for rational $P=P^\natural\cn$ the recursion computes $\cn_n=(\coef{z^n}R^{(n)})_{00}(0)\in\bbQ$.
\end{proof}

\begin{lem}[The $\bbZ/5$-degree of the gauge equation]
\label{lem:gau-degree}
For $P=\sum_kz^kP_k$ with $P_k\in\mathfrak{gl}_5(\Bbbk\lser u)$ and $\delta\in\bbZ/5$ put $P^{[\delta]}:=\sum_kz^k\sum_{(a,b):\,k+b-a\equiv\delta}(P_k)_{ab}E_{ab}$, so that $P=\sum_\delta P^{[\delta]}$, and $P$ is graded if and only if $P=P^{[0]}$. Let $\operatorname{ge}(P):=zP'-AP+P\tilde A$. Then $\operatorname{ge}(P^{[\delta]})=\operatorname{ge}(P)^{[\delta+1]}$; in particular every degree part of a solution of the gauge equation is a solution.
\end{lem}

\begin{proof}
Give $z^kf(u)E_{ab}$ the degree $k+b-a$. Then $zP'$ raises it by $1$, and so do $A\,(fz^kE_{cd})=b_{c-1}fz^kE_{c-1,d}$ and $(fz^kE_{cd})\tilde A=\tilde b_dfz^kE_{c,d+1}$; so $\operatorname{ge}$ maps degree $\delta$ to degree $\delta+1$.
\end{proof}

\begin{prop}[Graded gauges]
\label{prop:gau-graded}
Let $A$, $\tilde A$ have the one-pole shape and satisfy (EP).
\begin{enumerate}
\item Graded regular gauges exist; they can be chosen rational if all $b_a$, $\tilde b_a$ are rational.
\item If $P$ is a graded regular gauge, the graded regular gauges are exactly the $P\cn(z)$ with $\cn\in\Bbbk\fps{z^5}^\times$.
\item For a graded regular gauge $P$, $P_{11}(0,z)=\sum_{m\geq0}(P_{5m})_{11}(0)\,z^{5m}$ is a unit of $\Bbbk\fps{z^5}$.
\item If $\zeta\in\Bbbk$, then $\operatorname{Rot}_iP=\sum_\delta\zeta^{-i\delta}P^{[\delta]}$; so $P$ is graded if and only if $\operatorname{Rot}_1P=P$, if and only if $\operatorname{Rot}_iP=P$ for all $i$, and $\frac15\sum_i\operatorname{Rot}_iP=P^{[0]}$. Moreover $A\,\zeta^{i(\mathrm{gr}+1)}=\zeta^i\zeta^{i(\mathrm{gr}+1)}A$, and the same for $\tilde A$; $\operatorname{Rot}_i$ maps solutions of $z\de_uY=AY$ to solutions, and $\operatorname{Rot}_i(PY)=(\operatorname{Rot}_iP)(\operatorname{Rot}_iY)$, $\operatorname{Rot}_i(Y\mathsf d(z))=\operatorname{Rot}_i(Y)\,\mathsf d(\zeta^{-i}z)$ for scalar series $\mathsf d$.
\end{enumerate}
\end{prop}

\begin{proof}
(i) By \cref{lem:gau-degree}, $(P^\natural)^{[0]}$ solves the gauge equation, and its $z^0$-term is the diagonal $P^\natural_0$; rationality is \cref{thm:gau-U}~(d). (ii) Regular gauges are $P\cn$, $\cn\in\Bbbk\fps z^\times$ (\cref{thm:gau-U}~(b)); with $\cn^{(r)}:=\sum_{n\equiv r}\cn_nz^n$, $(P\cn)^{[r]}=P\cn^{(r)}$, and $\cn\mapsto P\cn$ is injective, so $P\cn$ is graded if and only if $\cn\in\Bbbk\fps{z^5}$. (iii) For graded $P$, $(P_k)_{11}=0$ unless $5\mid k$, and $(P_0)_{11}(0)=\cn(0)\mathfrak d(0)d^U_1(0)\neq0$. (iv) The entry $(a,b)$ of $\coef{z^k}\operatorname{Rot}_iP$ is $\zeta^{-i(k+b-a)}(P_k)_{ab}$; the equivalences follow from $\sum_i\zeta^{-i\delta}=5\delta_{\delta,0}$. The entry $(a,a+1)$ of $\zeta^{i(\mathrm{gr}+1)}A\zeta^{-i(\mathrm{gr}+1)}$ is $\zeta^{-i}b_a$, so for a solution $Y$, $z\de_u(\operatorname{Rot}_iY)=\zeta^i\zeta^{i(\mathrm{gr}+1)}A\,Y(u,\zeta^{-i}z)=A\operatorname{Rot}_iY$. The product rules are immediate.
\end{proof}

\input{sections/s03_gauge_5}



